%%%%%%%%%%%%%%%%%%%%%%%%%%%%%%%%%%%%%%%%%%%%%%%%%%%%%%%%%%%%%%%%%%%%%%
%  第 3 个逻辑月 · 月度报告
%  主题：形态 · 物理
%
%  说明：组长对此月只给出「形态」「物理」两个关键词，
%        具体所指【待确认3】。本报告按「形态定型 + 物理基础可信度」
%        的理解组织，相关判断待组长确认。
%  编译：xelatex 本文件（跑两遍）
%%%%%%%%%%%%%%%%%%%%%%%%%%%%%%%%%%%%%%%%%%%%%%%%%%%%%%%%%%%%%%%%%%%%%%
\documentclass[stage2,fontset=fandol]{ydzjreport}

\stagetwosetup
  {第 3 个逻辑月 · 形态与物理}%
  {冯旭东、陈思文、张昊洋}%
  {复合微通道分区冷却}%
  {【待确认】逻辑月 3 对应的自然月}%
  {形态定型；单柱基准；物性复算；网格无关性}%
  {把多柱整体对比收敛到单柱基准，剥离单参数影响；
   并推进新物性 UDF 的挂接与代表工况复算，补齐物理基础。}%

\begin{document}
\makestagetwohead

\begin{notebox}
\textbf{本月关键词说明}：组长提纲中本月只有「形态」「物理」两个词。
本报告按以下理解组织，\textbf{具体所指待确认}：
\begin{itemize}
  \item \textbf{形态}——把上月候选结构\textbf{收敛定型}，并回答"为什么是这个形态"；
  \item \textbf{物理}——\textbf{物理基础的可信度}：物性、流动模型、网格收敛与流场机理。
\end{itemize}
【待确认3】「形态」与「物理」的确切所指，以及二者是并行还是先后。
\end{notebox}

% ==================================================================
\section{本月研究任务}

\subsection{任务目标}

\begin{enumerate}
  \item 建立\textbf{单个扰流柱的基准模型}，逐一控制变量考察厚度、弯度、攻角的独立影响；
  \item 完成新水物性 UDF 的\textbf{挂接与自检}，并对代表工况\textbf{复算}；
  \item 收尾\textbf{网格无关性}检验，明确后续建模的网格基准；
  \item 开展流场与温度场的\textbf{机理分析}，解释已观察到的结构差异。
\end{enumerate}

\subsection{完成情况}

\begin{table}[H]
\centering
\caption{本月任务完成情况}
\label{tab:tasks}
\begin{tabularx}{\textwidth}{@{}Y c Y@{}}
\toprule
\rowcolor{ydzjlight} 任务 & 状态 & 说明 \\
\midrule
网格无关性检验 & 已完成 & 约 150 万网格作为后续划分基准 \\
新水物性 UDF 编写与编译 & 已完成 & 已识别原函数误差并完成编译 \\
单柱基准模型 & 未开始 & 上月已提出方向，本月尚未建立 \\
新 UDF 挂接与代表工况复算 & 未开始 & 编译已完成，挂接与复算待推进 \\
流场与温度场机理分析 & 部分完成 & 已定位高弯度翼型的尾缘流动死区 \\
\bottomrule
\end{tabularx}
\end{table}

% ==================================================================
\section{研究进展}

\subsection{已完成的工作}

\textbf{（一）网格无关性检验}

对直线、正弦、NACA2416 与 gao2.1 四种流道分别开展网格无关性检验，
比较不同网格密度下的计算结果，发现\textbf{网格数量达到约 150 万及以上时
表现较为理想}。

需要说明：该规模是本轮检验形成的\textbf{基准}，
\textbf{不表示}达到该数量即可自动满足所有模型的精度要求。
建模时以这一规模为起点，并在窄隙、尖角与流固界面局部加密；
几何或工况发生变化时仍需复核网格敏感性。

\begin{table}[H]
\centering
\caption{网格无关性检验结论}
\label{tab:mesh}
\begin{tabularx}{\textwidth}{@{}Y Y Y@{}}
\toprule
\rowcolor{ydzjlight} 流道 & 检验内容 & 结论 \\
\midrule
直线 & 多组网格密度对比 & 约 150 万及以上结果变化幅度很小 \\
正弦 & 多组网格密度对比 & 同上 \\
NACA2416 & 多组网格密度对比 & 同上 \\
gao2.1 & 多组网格密度对比 & 同上 \\
\bottomrule
\end{tabularx}
\end{table}

【缺数据：各结构的网格规模、关键物理量及相邻网格结果变化的归档表】

\textbf{（二）水物性 UDF 的修正}

上月核查发现原水物性函数存在误差。本月完成\textbf{新物性 UDF 的编写与编译}：
采用 101325 Pa 下 REFPROP WATER.FLD v9.1 液态水数据拟合，
提供随温度变化的密度、黏度、导热系数与比热；
\textbf{比热与显焓采用同一多项式及其积分}，保证二者定义一致。
源码对流体温度 280 至 360 K 进行范围检查，并提供 300 K 与 330 K 的物性自检输出。

工作由"发现误差"推进到"完成新函数编译"，
后续还需完成\textbf{挂接、物性校验与代表工况复算}。

\begin{table}[H]
\centering
\caption{原水物性在 298.15 K 校验点的偏差}
\label{tab:udf}
\begin{tabularx}{\textwidth}{@{}Y c Y@{}}
\toprule
\rowcolor{ydzjlight} 物性 & 相对参考值偏差 & 影响 \\
\midrule
密度 & 偏低约 60.8\% & 直接影响流量与压降 \\
比热 & 偏低约 10.5\% & 影响温升与热量收支 \\
黏度 & 偏低约 26.7\% & 影响雷诺数与流动阻力 \\
\bottomrule
\end{tabularx}
\end{table}

\textbf{（三）机理分析的初步结果}

针对上月发现的"高弯度翼型 PEC 反而回落"这一现象，
结合流场分析定位了原因：\textbf{高弯度翼型尾缘附近形成流动死区}，
分离使有效换热面积减小，抵消了弯度带来的增益。

\subsection{与预期的差异}

\textbf{一、物性误差的幅度超出预期。}
原本预期物性函数只是精度略有不足，
实测在 298.15 K 校验点密度偏低约 60.8\%，
\textbf{这不是精度问题，而是量级错误}。
这意味此前所有温度与压差结果的绝对值都不可信，
结构之间的排序也可能受影响。

\textbf{二、网格基准的可迁移性不能想当然。}
原本认为"加密到一定程度即可"，
检验后确认约 150 万是\textbf{本轮检验形成的参考规模}，
几何或工况变化时仍需复核，不能直接套用。

\subsection{尚存的疑问}

\begin{enumerate}
  \item 【待确认3】「形态」是否即指把上月候选收敛定型？「物理」指物性还是物理场机理？
  \item 修正物性之后，上月的形态排序（直线优于正弦、NACA2416 低阻）是否保持？
        需要\textbf{重新计算}才能判断，不能只重算后处理指标。
  \item 单柱模型应沿用哪一套网格与边界条件？
        若与整体流道不一致，单柱结论无法回推到整体。
\end{enumerate}

% ==================================================================
\section{问题与对策}

\subsection{遇到的主要问题}

\begin{enumerate}
  \item \textbf{新物性尚未挂接，现有结论仍基于旧物性。}
        修正工作只完成到"编译"，未完成"挂接与复算"，
        因此上月所有结论仍处在待验证状态。
  \item \textbf{单柱基准模型未建立。}
        上月已确认"多柱干扰无法解耦"，但没有单柱模型，
        参数影响无法定量剥离。
  \item \textbf{「物理」的工作边界不明确。}
        无法判断本月应做到"物性修正"为止，还是要一并完成流场机理的系统分析。
\end{enumerate}

\subsection{已尝试的方案}

\begin{enumerate}
  \item 针对问题 1，已完成新 UDF 的编写、编译与温度范围检查；
        挂接与复算需要在下月继续。
  \item 针对问题 2，已从流场分析入手定位现象原因（尾缘流动死区），
        属于定性解释；定量剥离仍需单柱模型。
  \item 针对问题 3，本报告按"物性 + 机理"两者都做的理解组织，
        待组长确认后可调整范围。
\end{enumerate}

\subsection{下一步对策}

\begin{enumerate}
  \item 挂接新 UDF，完成物性自检与代表工况复算，确认排序是否变化；
  \item 建立单柱基准模型，采用与整体流道一致的网格基准与边界条件；
  \item 在单柱模型上做单参数扫掠，得到厚度、弯度、攻角的独立影响曲线。
\end{enumerate}

% ==================================================================
\section{阶段性结论}

\subsection{可以确定的}

\begin{enumerate}
  \item 四种流道均已通过网格无关性检验，形成约 150 万网格的后续划分基准。
  \item 原水物性函数存在量级误差（密度偏低约 60.8\%），已识别并完成新 UDF 编译。
  \item 高弯度翼型 PEC 回落的原因已定位为\textbf{尾缘流动死区}导致有效换热面积减小。
  \item \textbf{物理基础必须先于结构比较完成}——
        在物性不可信的情况下比较结构，等于在错误的基准上排序。
\end{enumerate}

\subsection{尚不能确定的}

\begin{enumerate}
  \item \textbf{物性修正对结论的实际影响幅度。}
        缺少的证据是：挂接新 UDF 后代表工况的复算结果。
  \item \textbf{单柱条件下三个翼型参数的独立影响权重。}
        缺少的证据是：单柱基准模型上的单参数扫掠数据。
  \item \textbf{「形态」与「物理」的最终工作范围。}
        取决于组长对本月关键词的明确。
\end{enumerate}

% ==================================================================
\section{下月计划}

\begin{table}[H]
\centering
\caption{下月工作计划}
\label{tab:plan}
\begin{tabularx}{\textwidth}{@{}Y Y c@{}}
\toprule
\rowcolor{ydzjlight} 事项 & 预期产出 & 时间 \\
\midrule
挂接新 UDF，完成物性自检与代表工况复算 & 复算结果与排序对比 & 第 1 周 \\
建立单柱基准模型 & 单柱模型与基准数据 & 第 1--2 周 \\
单参数扫掠（厚度、弯度、攻角） & 参数影响特性曲线 & 第 2--4 周 \\
小块连接段设计（为下一逻辑月铺垫） & 连接段形态方案 & 第 3--4 周 \\
\bottomrule
\end{tabularx}
\end{table}

% ==================================================================
\begin{thebibliography}{9}
\bibitem{udf} 项目组. water\_properties\_refprop\_fit.c：液态水物性拟合与 Fluent UDF 源码. 项目计算资料, 2026.
\bibitem{iapws} IAPWS. SR6-08(2011)：0.1 MPa 下液态水物性补充发布. 2011.
\bibitem{monitor} Ansys. Fluent 用户指南：收敛监测. 2025 R1.
\end{thebibliography}

\end{document}
